\chapter{Performance}
\label{ch:perf}

Every figure in this chapter was produced by the test applications tc10, tc11 and tc12
of nornCore, run by \code{tools/measure.sh} on the Arduino Mega 2560 (16\,MHz) and the
ESP32 board (240\,MHz) and written to \code{doc/measurements/\measLabel.md} on
\measDate. The graphs are generated from that summary; a new measurement regenerates
them (\code{make data} in nornDocs). The same script records the ESP32 in QEMU, which is
not cycle accurate and is omitted here except where stack use is concerned, for which it
is exact.

\section{Method}

\paragraph{tc10} measures the kernel in CPU cycles with a cycle counter (a free running
timer at $F_{CPU}$ on the Mega, \code{CCOUNT} on the ESP32) and a probe interrupt whose
period is dithered so that it lands at every phase of the tick and of the other
activity. Everything in the test runs at once: the switch benchmark, the probe at about
770\,Hz, the tick, and the default configuration (stack fill, run time statistics,
assertions, software timers). Three kinds of figures come out:

\begin{itemize}
\item \textbf{API cost}: one call pair (give and take, lock and unlock, send and
  receive) on a single task with nobody waiting, so no switch happens. The minimum of
  64 repetitions: the cost of the critical section and the bookkeeping.
\item \textbf{Switch}: a low task signals a high one, which blocks again at once. Per
  switch, including the signalling call, minimum and average over 500 round trips. This
  is the number that matters for a design: how much a wake-up costs.
\item \textbf{Interrupt}: the entry latency from the hardware event to the first
  instruction of the \code{RTOS\_ISR} body, and the latency from a give in the ISR body
  to the first instruction of the woken task (``ISR to task''), minimum, average and
  maximum over a thousand probes landing at random phases.
\end{itemize}

simavr is cycle exact and gives the same numbers as the Mega board (the ISR to task
minimum is one lucky probe and may differ by a few cycles between the two), so a
simulator run of tc10 is the arbiter for an A/B comparison of two kernel variants.

\paragraph{tc11} measures the stack a task needs for the kernel's and the port's sake:
tasks that do one thing each (only block, be woken from an interrupt, be interrupted
while running, run the benchmark script) report their high water marks. \textbf{tc12}
repeats this at the bottom of a call chain 0, 4, 8, 16 and 32 frames deep, in four
phases.

\section{The cost of a call}

\begin{figure}[htbp]
\centering
\begin{tikzpicture}
\begin{axis}[nornbar, ylabel={cycles, one call pair},
  xticklabels from table={\apitable}{call}, legend entries={Mega 2560, ESP32}]
\addplot[fill=nornblue!70, draw=nornblue] table[x expr=\coordindex, y=mega2560] {\apitable};
\addplot[fill=nornred!60, draw=nornred] table[x expr=\coordindex, y=esp32] {\apitable};
\end{axis}
\end{tikzpicture}
\caption{API cost in cycles for one call pair without a switch (tc10, \measDate). The
AVR is an 8 bit machine; its cycle counts are higher but its cycle is 15 times longer,
so a semaphore pair costs 9\,$\mu$s on the Mega and 0.47\,$\mu$s on the ESP32.
\code{getStackFree} (1694 / 6438 cycles, the scan of a whole stack) is left out of the
graph.}
\label{fig:api}
\end{figure}

Figure~\ref{fig:api} ranks the primitives. The semaphore and the event flags are the
cheapest signals, the notification follows, the mutex costs a priority change on top of
the bookkeeping, the queue copies the item twice, and the pool pays for the pointer
alignment through \code{memcpy} on the AVR. \code{getCpuLoad} on the AVR is a 32 bit
division; \code{getTimeUs} on the ESP32 is a division of the cycle counter. ``Ceiling
held'' repeats the semaphore pair while the running task holds a mutex with a ceiling
above its priority: the scheduler then has to consider the boosted mask, and the
difference (8 cycles on the Mega, 3 on the ESP32) is the whole price of the ceiling
protocol on the hot path.

\section{The context switch}

\begin{figure}[htbp]
\centering
\begin{tikzpicture}
\begin{axis}[nornbar, ylabel={cycles per switch}, bar width=8pt,
  xticklabels from table={\switchtable}{mechanism},
  legend entries={Mega min, Mega avg, ESP32 min, ESP32 avg}, legend columns=2]
\addplot[fill=nornblue!80, draw=nornblue] table[x expr=\coordindex, y=megamin] {\switchtable};
\addplot[fill=nornblue!35, draw=nornblue] table[x expr=\coordindex, y=megaavg] {\switchtable};
\addplot[fill=nornred!70, draw=nornred] table[x expr=\coordindex, y=espmin] {\switchtable};
\addplot[fill=nornred!30, draw=nornred] table[x expr=\coordindex, y=espavg] {\switchtable};
\end{axis}
\end{tikzpicture}
\caption{Cycles per context switch including the signalling call, by mechanism
(tc10, \measDate): 27.4\,$\mu$s on the Mega and 1.70\,$\mu$s on the ESP32 for the
semaphore. The average exceeds the minimum on the Mega because the tick and the probe
interrupt land inside some of the 500 round trips.}
\label{fig:switch}
\end{figure}

Where the cycles go on the AVR, per switch: about 105 in the context switch itself (19
registers each way), about 45 in the run time statistics, about 20 in the assertions,
the rest in the scheduler, the wait list and the primitive. With statistics and
assertions off the switch is 322 cycles (20\,$\mu$s); table~\ref{tab:threecfg} has the
three configurations. On the ESP32 a switch spills every live register window and the
blocking path costs what a whole AVR stack does; with statistics off it is 361 cycles,
with assertions off as well 351 (1.46\,$\mu$s).

\section{Interrupt latency}

\begin{figure}[htbp]
\centering
\begin{tikzpicture}
\begin{axis}[nornbar, ylabel={cycles, ISR body to task}, bar width=8pt, enlarge x limits=0.25,
  xticklabels from table={\isrtable}{path},
  legend entries={Mega min, Mega avg, Mega max, ESP32 min, ESP32 avg, ESP32 max}, legend columns=3]
\addplot[fill=nornblue!90, draw=nornblue] table[x expr=\coordindex, y=megamin] {\isrtable};
\addplot[fill=nornblue!55, draw=nornblue] table[x expr=\coordindex, y=megaavg] {\isrtable};
\addplot[fill=nornblue!20, draw=nornblue] table[x expr=\coordindex, y=megamax] {\isrtable};
\addplot[fill=nornred!90, draw=nornred] table[x expr=\coordindex, y=espmin] {\isrtable};
\addplot[fill=nornred!55, draw=nornred] table[x expr=\coordindex, y=espavg] {\isrtable};
\addplot[fill=nornred!20, draw=nornred] table[x expr=\coordindex, y=espmax] {\isrtable};
\end{axis}
\end{tikzpicture}
\caption{Latency from a give in the ISR body to the first instruction of the woken task,
minimum, average and maximum over 1000 probes at random phases (tc10, \measDate). The
maximum is the longest interrupt-off window: a context switch that the probe lands
behind.}
\label{fig:isr}
\end{figure}

\begin{table}[htbp]
\centering\small
\begin{tabular}{@{}lrrrr@{}}
\toprule
 & \multicolumn{2}{c}{Mega 2560} & \multicolumn{2}{c}{ESP32} \\
 & cycles & time & cycles & time \\
\midrule
interrupt entry, min / avg / max         & 50 / 65 / 548 & 3.1 / 4.1 / 34 $\mu$s & -- / 35 / 106 & -- / 0.9 / 2.7 $\mu$s \\
ISR to task (semaphore), min / max       & 471 / 826 & 29.4 / 51.6 $\mu$s & 418 / 765 & 1.74 / 3.19 $\mu$s \\
hardware event to task, best case        & & 33.5 $\mu$s & & 1.85 $\mu$s \\
\bottomrule
\end{tabular}
\caption{Interrupt figures (tc10, \measDate). The ESP32 entry minimum reads 0 because the
probe timer's 40\,MHz counter is latched at the alarm. The entry cost on the ESP32 is
the same 850\,ns whether the interrupted task is 0 or 32 call frames deep (tc12): the
window spill at entry depends on the live windows, at most about seven, not on the
depth.}
\label{tab:isr}
\end{table}

The AVR entry maximum of 34\,$\mu$s and the ISR to task maximum of 52\,$\mu$s are the
same phenomenon: a 27\,$\mu$s switch with interrupts off that the probe lands behind,
plus a tick with a wake-up. Interrupts do not nest, so these maxima are also the worst
case for any other interrupt in the system.

\section{Memory}

\begin{figure}[htbp]
\centering
\begin{tikzpicture}
\begin{axis}[nornbar, ylabel={bytes}, bar width=8pt,
  xticklabels from table={\ramtable}{object}, legend entries={Mega 2560, ESP32}]
\addplot[fill=nornblue!70, draw=nornblue] table[x expr=\coordindex, y=mega2560] {\ramtable};
\addplot[fill=nornred!60, draw=nornred] table[x expr=\coordindex, y=esp32] {\ramtable};
\end{axis}
\end{tikzpicture}
\caption{RAM per kernel object in the default configuration (tc10, \measDate). The TCB
is per task; the pointer width (2 and 4 bytes) explains most of the difference.}
\label{fig:ram}
\end{figure}

\begin{table}[htbp]
\centering\small
\begin{tabular}{@{}lrrrrrr@{}}
\toprule
 & \multicolumn{3}{c}{image (tc01)} & \multicolumn{3}{c}{every kernel and port function} \\
Board & text & data & bss & text & data & bss \\
\midrule
uno       & 6852  & 170 & 833  & 9909  & 8  & 240 \\
mega2560  & 6998  & 170 & 833  & 9935  & 8  & 240 \\
esp32     & 10138 & 24  & 3920 & 12975 & 13 & 815 \\
\bottomrule
\end{tabular}
\caption{Code and data size in bytes (\measLabel). The image is tc01 with its default
configuration, after the linker dropped what it does not call; the right half is every
kernel and port object compiled, the upper bound an application can pull in. The
smallest running application, one task blinking the LED with everything off, is 2514
bytes of text on the UNO and 6089 on the ESP32.}
\label{tab:size}
\end{table}

\section{Stacks}

\begin{figure}[htbp]
\centering
\begin{tikzpicture}
\begin{axis}[nornbar, ylabel={bytes used}, bar width=10pt, enlarge x limits=0.2,
  xticklabels from table={\stackeleventable}{task}, legend entries={Mega 2560, ESP32}]
\addplot[fill=nornblue!70, draw=nornblue] table[x expr=\coordindex, y=mega2560] {\stackeleventable};
\addplot[fill=nornred!60, draw=nornred] table[x expr=\coordindex, y=esp32] {\stackeleventable};
\end{axis}
\end{tikzpicture}
\caption{Stack used by tasks that do one thing each (tc11, \measDate): \emph{block} only
blocks, \emph{wake} is woken from an interrupt, \emph{busy} is interrupted while
running, \emph{script} runs the test's own sequence. The ESP32 ``block'' figure is 192
when no interrupt lands on the task during its short runs and 300 when one does: size a
stack by the interrupted figure, which does not depend on luck.}
\label{fig:stack11}
\end{figure}

\begin{figure}[htbp]
\centering
\begin{tikzpicture}
\begin{axis}[nornline, width=0.49\linewidth, height=6cm, xlabel={call depth (frames)},
  ylabel={bytes used}, title={Mega 2560}, xticklabels from table={\stackmegatable}{depth},
  legend entries={A chain alone, B preempted, C blocking, D float preempted}, legend style={font=\scriptsize}]
\addplot[mark=*, gray] table[x expr=\coordindex, y=A] {\stackmegatable};
\addplot[mark=square*, nornblue] table[x expr=\coordindex, y=B] {\stackmegatable};
\addplot[mark=triangle*, nornblue!50] table[x expr=\coordindex, y=C] {\stackmegatable};
\addplot[mark=diamond*, nornred] table[x expr=\coordindex, y=D] {\stackmegatable};
\end{axis}
\end{tikzpicture}\hfill
\begin{tikzpicture}
\begin{axis}[nornline, width=0.49\linewidth, height=6cm, xlabel={call depth (frames)},
  ylabel={bytes used}, title={ESP32}, xticklabels from table={\stackesptable}{depth},
  legend entries={A chain alone, B preempted, C blocking, D float preempted}, legend style={font=\scriptsize}]
\addplot[mark=*, gray] table[x expr=\coordindex, y=A] {\stackesptable};
\addplot[mark=square*, nornblue] table[x expr=\coordindex, y=B] {\stackesptable};
\addplot[mark=triangle*, nornblue!50] table[x expr=\coordindex, y=C] {\stackesptable};
\addplot[mark=diamond*, nornred] table[x expr=\coordindex, y=D] {\stackesptable};
\end{axis}
\end{tikzpicture}
\caption{Stack used against call depth (tc12, \measDate): the chain alone, preempted at
the bottom with a switch out of the interrupt, blocking at the bottom, and floating point
at the bottom. The lines are parallel: the kernel's and the port's share is constant
(Mega 59 / 48 / 65 bytes, ESP32 368 / 300 / 368 including the task's two bottom
frames), so the rule for a stack is the task's own frames (12 bytes per level on the
Mega with a few locals, 32 per \code{call8} level on the ESP32) plus the
``interrupted'' figure.}
\label{fig:stack12}
\end{figure}

\section{The cost of the configuration}

Table~\ref{tab:threecfg} gives the switch for the three configurations: on the Mega the
run time statistics cost about 45 cycles per switch and the assertions about 20; on the
ESP32 the lean kernel saves 93 cycles over the default. \code{RTOS\_CFG\_STACK\_CHECK}
costs nothing on the hot path (the fill is at creation, the canary check a compare per
switch), the software timers one list check per tick, and the ceiling protocol 8 cycles
on the Mega only while a boosting mutex is held.

\section{History}

nornCore measures itself before and after every change to the kernel. Figure~\ref{fig:history}
shows the hot path and the code size across the steps since the baseline of
2026-11-05: the fixes of an independent audit (F1 to F11), the scheduler contract for
mutex ceilings (tie), the exclusion of the running task's own boost from the slow path
test (boost) and the task life cycle (lifecycle). Correctness bought about 13 cycles per
switch on the Mega (F2: the scheduler now honours every task's effective priority, not
only the running one's) and about 550 bytes of flash; the later steps are flat.

\begin{figure}[htbp]
\centering
\begin{tikzpicture}
\begin{axis}[nornline, height=5.5cm, ylabel={cycles}, xticklabels from table={\historytable}{step},
  legend entries={Mega switch, Mega ISR to task, ESP32 switch, ESP32 ISR to task}, legend columns=2,
  legend style={at={(0.5,-0.42)}, anchor=north}]
\addplot[mark=*, nornblue] table[x=n, y=megaswitch] {\historytable};
\addplot[mark=square*, nornblue!50] table[x=n, y=megaisr] {\historytable};
\addplot[mark=*, nornred] table[x=n, y=espswitch] {\historytable};
\addplot[mark=square*, nornred!50] table[x=n, y=espisr] {\historytable};
\end{axis}
\end{tikzpicture}

\vspace{1em}
\begin{tikzpicture}
\begin{axis}[nornline, height=5.5cm, ylabel={text, bytes (tc01)}, xticklabels from table={\historytable}{step},
  legend entries={uno, mega2560, esp32}, legend columns=3, legend style={at={(0.5,-0.42)}, anchor=north}]
\addplot[mark=*, gray] table[x=n, y=unotext] {\historytable};
\addplot[mark=*, nornblue] table[x=n, y=megatext] {\historytable};
\addplot[mark=*, nornred] table[x=n, y=esptext] {\historytable};
\end{axis}
\end{tikzpicture}
\caption{The hot path (semaphore switch and ISR to task minimum) and the image size
across the measured steps in \code{doc/measurements/}.}
\label{fig:history}
\end{figure}

\section{The older timing test}

tc04 measures with a 4\,$\mu$s counter and runs on every board including the UNO; it is
coarser but confirms the picture: 24\,$\mu$s per switch and 24\,$\mu$s worst interrupt
latency on both AVR boards, 1.7\,$\mu$s and 2\,$\mu$s on the ESP32 at 240\,MHz (10.9 and
12\,$\mu$s at the 40\,MHz crystal clock before the PLL is enabled), with a periodic jitter
of 16\,$\mu$s on the AVR and none the counter can see on the ESP32.
